\documentclass[aps,pra,twocolumn,superscriptaddress,nofootinbib]{revtex4-2}

\usepackage{amsmath,amssymb,mathtools}
\usepackage{bbm}
\usepackage{amsthm}
\usepackage{booktabs}
\usepackage{array}
\usepackage{graphicx}
\usepackage[colorlinks=true,linkcolor=blue,citecolor=blue,urlcolor=blue]{hyperref}

\newtheorem{definition}{Definition}
\newtheorem{proposition}{Proposition}
\newtheorem{remark}{Remark}

\newcommand{\C}{\mathcal{C}}
\newcommand{\E}{\mathcal{E}}
\newcommand{\F}{\mathcal{F}}
\newcommand{\HH}{\mathcal{H}}
\newcommand{\OmegaSet}{\Omega}
\newcommand{\TVS}{\mathcal{S}_2}
\newcommand{\Ptwo}{\mathsf{P}_{2}}
\newcommand{\Pcol}{\mathsf{P}_{\mathrm{col}}}
\newcommand{\Cone}{\operatorname{Cone}}
\newcommand{\conv}{\operatorname{conv}}
\newcommand{\rank}{\operatorname{rank}}
\newcommand{\tr}{\operatorname{tr}}
\newcommand{\id}{\mathbbm{1}}
\newcommand{\reals}{\mathbb{R}}
\newcommand{\complex}{\mathbb{C}}
\newcommand{\diag}{\operatorname{diag}}
\newcommand{\YO}{\mathrm{YO}}

\begin{document}

\title{Combinatorial, Simplex, and Operator-Functional Contextuality in Quantum Logics}

\author{Karl Svozil}
\affiliation{Institute for Theoretical Physics, TU Wien, Wiedner Hauptstra{\ss}e~8--10/136, 1040 Vienna, Austria}

\date{\today}

\begin{abstract}
This paper separates three classicality tests that are often conflated in discussions of contextuality: combinatorial tests based on valuations, colorings, and partition-logic embeddability; simplex tests for operational probability tables; and operator-functional tests requiring product- and function-preserving values of observables.  The distinction is relative to the structure retained in the fragment---incidence data, probability data, or operator algebra.  Calibration examples and projector-cone computations show that these tests often coexist in familiar quantum examples, but none uniformly replaces the others.
\end{abstract}

\maketitle


\section{Introduction}

In Hilbert spaces of dimension at least three, quantum propositions cannot in general be assigned pre-existing, context-independent truth values compatible with the functional structure of sharp measurements~\cite{Gleason,specker-60,kamber65,ZirlSchl-65,Zierler1975,kochen1,cabello-96}.  In the language of quantum logic, this obstruction appears as the absence, scarcity, or degeneracy of admissible two-valued states on a pasted family of Boolean algebras, an orthomodular poset, or an orthogonality hypergraph.  At the probabilistic level, the same issue can be reformulated operationally: one asks whether a specified prepare-and-measure experiment admits a classical interpretation~\cite{froissart-81,Garg1984,pitowsky-86,Klyachko-2008,Yu-2012}.  Certain operator-valued configurations sharpen this tension into state-independent contradictions between quantum and classical predictions, provided one assumes simultaneous pre-existing values for individual counterfactual observables~\cite{ghz,mermin90b,svozil-2024-convert-pra-externalfigures}.  Generalized-noncontextual ontological models~\cite{Spekkens-04}, formalized through the simplex-embeddability criterion of Schmid, Selby, Wolfe, Kunjwal, and Spekkens~\cite{schmid-2021-simplex}, and implemented computationally by the linear program of Selby, Wolfe, Schmid, Sainz, and Rossi~\cite{selby-2024-lp}, provide a test for this operational notion of classicality.

The three criteria are related, but they are not the same test.  \emph{Combinatorial contextuality} tests an incidence structure: does the hypergraph or pasted Boolean structure admit enough admissible two-valued states, a partition-logic embedding~\cite{kochen1,svozil-2001-eua}, or a suitable chromatic coloring~\cite{svozil-2025-color}?  \emph{Simplex contextuality} tests an operational probability table: do the chosen preparations and effects admit a simplex embedding, possibly after adding a specified kind of noise?  \emph{Operator-functional contextuality} tests an operator representation: can values be assigned to observables while preserving the functional and product relations among commuting operators?  The aim of the paper is to keep these three questions distinct.

This taxonomy is meant to complement, not replace, existing classifications.  The graph-theoretic approach of Cabello, Severini, and Winter relates exclusivity graphs to invariants such as independence numbers and the Lov\'asz theta function~\cite{cabello-severini-winter-2014}.  The hypergraph approach of Ac\'in, Fritz, Leverrier, and Sainz formulates contextuality scenarios directly as compatibility/exclusivity structures with probabilistic models~\cite{acin-fritz-leverrier-sainz-2015}.  The sheaf-theoretic framework of Abramsky and Brandenburger classifies empirical models by the obstruction to gluing local sections, distinguishing probabilistic, logical, and strong forms of contextuality~\cite{abramsky-brandenburger-2011}.  The present distinction is cross-cutting: it asks which part of the Hilbert-space description is being retained before the classicality test is applied---incidence data, convex-operational probability data, or operator-algebraic product data.

The separation claims below should be read in this relative sense.  If a fragment is enlarged by adding unit effects, complements, coarse grainings, preparation equivalences, local factor measurements, or product constraints, then a phenomenon invisible at one level may become visible at another.  Thus the statement that the collective Greenberger--Horne--Zeilinger (GHZ) product measurement is simplex-trivial concerns only that closed single-context fragment; it does not deny that local GHZ statistics, or an unmasked Peres--Mermin fragment, can violate a full operational noncontextuality test.

The operator-functional layer is included because it behaves differently from both the projector-hypergraph and probability-table layers.  In GHZ- or Mermin-type parity arguments, the contradiction is not always located at the level of a single resolved Boolean context.  For the closed collective GHZ measurement, the commuting global product observables have a common eigenbasis, and their joint spectral projectors form one Boolean algebra.  The contradiction appears only when one additionally requires the assigned value of a global product, such as $X_1Y_2Y_3$, to factor into pre-existing local values of the observables $X_i$ and $Y_i$.  The classical product of the four assigned global values is then forced to be $+1$, whereas the corresponding operator product is $-\id$.  Thus the relevant issues are the assignment of values to individual observables and the product-preservation assumption, rather than the incidence structure of the resulting joint spectral projectors.  This does not place operator-valued parity proofs outside simplex analysis in general: the simplex verdict depends on the chosen operational fragment.  A closed collective single-context GHZ fragment is simplex-trivial, whereas an enlarged or spectrally unmasked fragment, including the relevant local factors, refinements, and operational equivalences, may become simplex-testable.  Peres--Mermin provides the standard case in which the operator parity proof descends, after unmasking, to a genuine multi-context projector configuration~\cite{mermin90b,svozil-2024-convert-pra-externalfigures}.

This terminology separates questions that are often conflated.  A hypergraph may be rich in two-valued states and even be representable by a partition logic, hence by a generalized urn or finite automaton~\cite{svozil-2001-eua}, while still not being a faithful quantum orthogonality structure.  Conversely, a quantumly coordinatized projector fragment may be simplex contextual because its Born probabilities cannot be mimicked by a classical simplex, even though the corresponding hypergraph is not a Kochen--Specker (KS) set in the strongest valuation sense.  Finally, an operator-valued parity argument may survive even when the extracted joint spectral context is Boolean.  The bare logic of contexts therefore does not by itself determine the probability theory or the operator-functional constraints.

Within the combinatorial layer, chromatic contextuality has recently been emphasized as a further global test~\cite{svozil-2021-chroma,svozil-2025-color}.  Given a $d$-uniform orthogonality hypergraph, one asks whether the vertices can be colored with $d$ colors so that every context receives all colors exactly once~\cite{svozil-2025-color}.  This is equivalent to asking for a global, context-independent assignment of a fixed nondegenerate spectrum to all contexts.  Failure of such colorability is a strong global obstruction, but it is neither identical to the absence of two-valued states nor identical to simplex nonembeddability.  It is also distinct from geometric coordinatizability.  For example, the completed Yu--Oh 25-ray configuration~\cite{Yu-2012} is faithfully coordinatized in $\mathbb R^3$ although it is not strongly 3-colorable; conversely, Greechie's $G_{32}$~\cite{greechie:71,Holland1975,Bennett-MC-1970,Greechie1974,Greechie-Suppes1976} has a rich supply of two-valued states but fails coordinatizability for projective-geometric reasons.  A companion analysis gives the detailed proof of this chromatic/coordinatization separation~\cite{svozil-2026-chromaticcoordinatization}; here it is used only as orientation for the taxonomy.

The paper proceeds as follows.  Section~\ref{sec:three-layers} defines the three-layer organization: combinatorial, simplex, and operator-functional contextuality.  Section~\ref{sec:simplex-test} recalls the linear-programming (LP) test for simplex embeddability.  Section~\ref{sec:sandwich} develops the geometric \emph{simplicial-sandwich} interpretation of simplex nonclassicality.  Section~\ref{sec:operator-valued} explains how operator-valued simplex tests depend on whether operators are treated as spectral-outcome data or as algebraic quantities subject to product rules.  Section~\ref{sec:two-contexts} explains why two contexts, or weakly pasted contexts, need not by themselves separate vector probabilities from classical convex probabilities.  Section~\ref{sec:calibration} discusses the firefly logic $L_{12}$ and the completed pentagon, including partition-logical, generalized-urn, and finite-automaton realizations.  The subsequent sections review meager two-valued-state spaces and chromatic colorability.

\section{Three layers of contextuality}
\label{sec:three-layers}

\subsection{Combinatorial contextuality}

Combinatorial contextuality is a property of the pasted context structure.  Its input is an orthogonality hypergraph $H=(V,\C)$ whose vertices are atoms and whose hyperedges are contexts.  The basic question is whether there exist global context-independent assignments compatible with the local Boolean structure.  The most familiar assignments are two-valued states
\begin{equation}
        s:V\to\{0,1\},
        \qquad \sum_{v\in C}s(v)=1\quad(C\in\C).
\end{equation}
Their absence gives the KS obstruction.  Their scarcity gives the meager cases considered below: nonunital, nonseparating, and nonfull two-valued-state spaces.  Their abundance, especially a separating family, permits a concrete set representation and hence a partition-logic model~\cite{svozil-2001-eua}.

Chromatic contextuality~\cite{svozil-2025-color} also belongs to the combinatorial level, but it asks for more than the existence of valuations.  A strong $d$-coloring of a $d$-uniform hypergraph decomposes the constant-one assignment into $d$ two-valued states, one per color.  Thus chromatic contextuality is a spectral-labeling obstruction: there is no way to use one fixed set of $d$ nondegenerate outcome labels globally across all contexts.  This is a discrete global obstruction, not a probability-table obstruction.

Partition logics sit at the opposite end of the same combinatorial level.  Their existence demonstrates that non-Boolean pasting and complementarity can also be classically represented.  A partition logic is obtained by taking a family of partitions of a finite set and identifying equal parts across different partitions.  Its dispersion-free states are the underlying points of the set.  In Wright-style generalized urns, the same construction is realized by ball types and color filters; in Moore or Mealy automata, it is realized by initial states and input-dependent output partitions.  Such models implement complementarity without value indefiniteness~\cite{schaller-92,svozil-2018-b}.

\subsection{Simplex contextuality}

Simplex contextuality is an operational convex-geometric notion.  Its input is not just a hypergraph but a finite prepare-and-measure fragment
\begin{equation}
        (\OmegaSet,\E,p(e|\rho)),
\end{equation}
where $\OmegaSet$ is a set of preparations, $\E$ is a set of effects, and $p(e|\rho)$ is the observed probability table.  In the quantum cases below, $p(e|\rho)=\tr(e\rho)$ and the preparations and effects are rank-one projectors associated with the listed rays.

A fragment is simplex embeddable if its statistics factor through a classical simplex.  Equivalently, the fragment has a generalized-noncontextual explanation in the sense of generalized probabilistic theories (GPTs).  If no such embedding exists, or if it exists only after adding noise, the fragment is simplex contextual.  This property depends on which preparations and effects are included.  Enlarging a ray set can change the generated cones and therefore the simplex robustness, even if the additional projectors are merely an orthogonal completion of a previously studied hypergraph.

The two levels can coincide, but need not.  A KS set such as the Cabello 18--9 configuration is nonclassical at the valuation level and, for the natural projector fragment, simplex contextual as well.  But a partition logic may be fully classical combinatorially while having a quantum double with different Born probabilities.

Conversely, a restricted operational fragment may fail to reveal a defect already present in the underlying hypergraph.  For example, nonunitality is exposed only if the offending atom can actually be prepared or measured; nonseparability is exposed only if the fragment contains preparations that statistically distinguish the identified atoms; and a true-implies-false set of two-valued states (TIFS) or nonfullness constraint~\cite{2018-minimalYIYS} is exposed only if the relevant endpoint preparations and effects are included.  Without such operational witnesses, the simplex test may remain feasible even though the incidence structure has a combinatorial defect.


\subsection{Operator-functional contextuality}

Operator-functional contextuality concerns observables as algebraic quantities, not merely as names or labels for their spectral projectors.  A classical value assignment to observables is then required to preserve functional relations, for example
\begin{equation}
        v(f(A))=f(v(A))
\end{equation}
and, for mutually commuting observables,
\begin{equation}
        v(AB)=v(A)v(B).
        \label{eq:product-rule}
\end{equation}
This requirement is stronger than assigning value one to a single atom in each joint spectral basis.  It also differs from a simplex test on probabilities, because it constrains the deterministic vertices by algebraic product identities.

This distinction is decisive for GHZ-type arguments.  The Mermin--GHZ operators
\begin{equation}
X_1X_2X_3,\quad X_1Y_2Y_3,\quad Y_1X_2Y_3,\quad Y_1Y_2X_3
\end{equation}
commute and therefore possess a common eigenbasis.  That eigensystem is one orthonormal basis in $\complex^8$; as a projector logic it is a single Boolean block with eight atoms and therefore has eight dispersion-free states~\cite{svozil-2024-convert-pra-externalfigures}.  There is no KS contradiction at the level of that isolated context.  The contradiction enters only when the value of each global product observable is assumed to factor into values of the local Pauli components $X_i$ and $Y_i$.  Then each local factor appears twice in the product of the four assigned values, forcing $+1$ classically, whereas the corresponding operator product is $-\id$.  The failure is therefore not the absence of a valuation on one context, but the impossibility of a product-preserving valuation on the operator algebra being implicitly invoked.

The Peres--Mermin square~\cite{peres111,mermin90b,peres-91}
is different but related.  At the operator level it is a parity proof.
When the masked contexts are extracted by simultaneous diagonalization or by matrix pencils,
the operator argument can be transcribed into a projective KS structure,
in particular its associated 24--24 (vector--context) completion~\cite{Pavicic-2005,svozil-2024-convert-pra-externalfigures}  containing the Cabello--Estebaranz--Garc\'ia-Alcaine 18--9 set~\cite{cabello-96,svozil-2024-convert-pra-externalfigures}.  Thus some operator-functional contradictions descend to combinatorial projector contextuality after unmasking, while others, such as the single-context GHZ eigensystem, do not.  This is why operator-functional contextuality deserves to be kept as a third layer rather than absorbed into either colorability or simplex embeddability.

The word ``single-context'' is therefore a statement about one particular operational packaging, not about the full GHZ reasoning.  The collective measurement of the four global products is a single Boolean context.  The local GHZ experiment, however, invokes four incompatible local settings and the counterfactual identification of local factors across them.  In the first packaging the simplex and two-valued-state tests are trivial; in the second packaging the nonclassicality is absolute and rests on locality, noncontextuality, and product preservation. Equivalently, the collective product observable and its local-factor implementation should not be identified without further assumptions.  The former reports a global spectral outcome; the latter refines that outcome into local Pauli outcomes and then imposes a product rule on their assigned values.

\subsection{Probability theories on the same pasting}

The probability level is itself not exhausted by the binary alternative ``classical or quantum.''  On a fixed exclusivity structure one may consider at least three types of context-independent additive weights.  First, classical probabilities are convex mixtures of dispersion-free states; they form a polytope.  Second, Born probabilities arise from a faithful orthogonal representation and a state vector or density operator; graph-theoretically, these are related to Lov\'asz-type theta-body constructions.  Third, there are exotic weights, such as Wright-type weights on pentagon-like logics~\cite{greechie-1974,wright:pent}, which satisfy additivity within contexts but are neither classical convex-hull points nor Hilbert-space Born probabilities.

Thus the incidence structure alone does not determine the probability calculus.  The same pasted logic may admit a partition model, a vector representation, and additional exotic additive weights.  This is the sense in which combinatorial contextuality and simplex contextuality are largely independent diagnostics.  The former classifies global assignments on the logic; the latter tests a particular operational probability model.


\section{The operational simplex test}
\label{sec:simplex-test}

Let $\OmegaSet=\{\rho_1,\ldots,\rho_N\}$ be a finite set of quantum states, and let
$\E=\{e_1,\ldots,e_M\}$ be a finite set of quantum effects.  The states may be
normalized or subnormalized.  Here subnormalized means
\begin{equation}
        \rho_i\ge 0,
        \qquad
        0\le \tr(\rho_i)\le 1.
\end{equation}
Thus one may write $\rho_i=p_i\widehat{\rho}_i$, where
$\widehat{\rho}_i$ is a normalized density operator and
$0\le p_i\le 1$ is a weight or success probability.  This convention is
natural for cone constructions, since the cones below are closed under
nonnegative rescaling.

Let
\begin{equation}
        \mathcal A=\operatorname{Herm}(\HH)
\end{equation}
denote the real vector space of Hermitian operators on $\HH$.  We do not
assume that the chosen coordinates on $\mathcal A$ are trace-orthonormal.
Instead, choose an arbitrary real basis
\begin{equation}
        B_1,\ldots,B_q
\end{equation}
of $\mathcal A$.  If
\begin{equation}
        A=\sum_{\alpha=1}^q a_\alpha B_\alpha,
        \qquad
        C=\sum_{\beta=1}^q c_\beta B_\beta,
\end{equation}
then the trace pairing is represented by the Gram matrix
\begin{equation}
        G_{\alpha\beta}=\tr(B_\alpha B_\beta).
        \label{eq:trace-gram-matrix}
\end{equation}
Thus, in these coordinates,
\begin{equation}
        \tr(AC)=a^T G c.
        \label{eq:trace-pairing-coordinate}
\end{equation}
Only in a trace-orthonormal basis would $G$ be the identity matrix.

Let
\begin{equation}
        S_\Omega=\operatorname{span}(\OmegaSet),
        \qquad
        S_E=\operatorname{span}(\E)
\end{equation}
be the accessible state and effect subspaces.  Choose arbitrary real bases of
$S_\Omega$ and $S_E$.  Let
\begin{equation}
        d_\Omega=\dim S_\Omega,
        \qquad
        d_E=\dim S_E.
\end{equation}
The inclusion maps
\begin{equation}
        I_\Omega:S_\Omega\hookrightarrow \mathcal A,
        \qquad
        I_E:S_E\hookrightarrow \mathcal A
\end{equation}
are the coordinate embeddings determined by these choices of bases.  Concretely,
$I_\Omega$ is a $q\times d_\Omega$ matrix whose columns are the ambient
coordinates, in the basis $B_1,\ldots,B_q$, of the chosen basis vectors of
$S_\Omega$.  Similarly, $I_E$ is a $q\times d_E$ matrix whose columns are the
ambient coordinates of the chosen basis vectors of $S_E$.

Therefore, if $x\in\reals^{d_\Omega}$ is the coordinate vector of an accessible
state and $y\in\reals^{d_E}$ is the coordinate vector of an accessible effect,
then their ambient coordinate vectors are $I_\Omega x$ and $I_E y$,
respectively.  The Born pairing between them is
\begin{equation}
        \tr(e_y\rho_x)
        =
        (I_E y)^T G (I_\Omega x)
        =
        y^T I_E^T G I_\Omega x.
        \label{eq:born-pairing-coordinate}
\end{equation}
Thus
\begin{equation}
        B_{E\Omega}=I_E^T G I_\Omega
        \label{eq:born-pairing-matrix}
\end{equation}
is the matrix of the accessible state--effect pairing in the chosen reduced
coordinates.

After choosing the reduced bases of $S_\Omega$ and $S_E$, let $V_\Omega$ and
$V_E$ denote the matrices whose columns are the coordinate vectors of the
$\rho_i$ and $e_j$, respectively.  The state and effect cones are first given
by their generator, or $V$-, representations:
\begin{align}
        \Cone[\OmegaSet]
        &=
        \left\{
        V_\Omega r : r\in\reals^N,\ r_i\ge 0
        \right\}                                                   \nonumber\\
        &=
        \left\{
        \sum_{i=1}^N r_i\rho_i : r_i\ge 0
        \right\}
        \subseteq S_\Omega,
        \label{eq:state-cone-vrep}
        \\
        \Cone[\E]
        &=
        \left\{
        V_E s : s\in\reals^M,\ s_j\ge 0
        \right\}                                                   \nonumber\\
        &=
        \left\{
        \sum_{j=1}^M s_j e_j : s_j\ge 0
        \right\}
        \subseteq S_E.
        \label{eq:effect-cone-vrep}
\end{align}

Because these cones are finitely generated, they are polyhedral.  Hence they
also admit half-space, or $H$-, representations.  That is, there are linear
functionals
\begin{equation}
        h^\Omega_1,\ldots,h^\Omega_{n_\Omega}\in S_\Omega^*,
        \qquad
        h^E_1,\ldots,h^E_{n_E}\in S_E^*,
\end{equation}
which may be chosen as facet-defining inequalities, such that
\begin{align}
        \Cone[\OmegaSet]
        &=
        \left\{
        x\in S_\Omega :
        h^\Omega_k(x)\ge 0
        \ \text{for } k=1,\ldots,n_\Omega
        \right\},
        \label{eq:state-cone-hrep}
        \\
        \Cone[\E]
        &=
        \left\{
        y\in S_E :
        h^E_\ell(y)\ge 0
        \ \text{for } \ell=1,\ldots,n_E
        \right\}.
        \label{eq:effect-cone-hrep}
\end{align}
Equivalently, collect these functionals as the rows of linear maps
\begin{equation}
        H_\Omega:S_\Omega\to\reals^{n_\Omega},
        \qquad
        H_E:S_E\to\reals^{n_E},
\end{equation}
so that
\begin{equation}
        H_\Omega x
        =
        \begin{pmatrix}
        h^\Omega_1(x)\\
        \vdots\\
        h^\Omega_{n_\Omega}(x)
        \end{pmatrix},
        \qquad
        H_E y
        =
        \begin{pmatrix}
        h^E_1(y)\\
        \vdots\\
        h^E_{n_E}(y)
        \end{pmatrix}.
\end{equation}
With this notation,
\begin{align}
        H_\Omega x\ge_{\mathrm e}0
        &\Longleftrightarrow
        x\in\Cone[\OmegaSet],
        \nonumber\\
        H_E y\ge_{\mathrm e}0
        &\Longleftrightarrow
        y\in\Cone[\E].
        \label{eq:facet-map-criterion}
\end{align}
Here $\ge_{\mathrm e}$ denotes entrywise nonnegativity: all components of the
corresponding vector must be nonnegative.  It is not the positive-semidefinite
order on Hermitian operators.

The simplex-embeddability linear program asks whether the Born-pairing matrix
\eqref{eq:born-pairing-matrix} factors through the facet descriptions of the
state and effect cones.  Equivalently, one asks whether there exists an
entrywise nonnegative matrix
\begin{equation}
        \sigma\in\reals^{n_E\times n_\Omega}
\end{equation}
such that
\begin{equation}
        I_E^T G I_\Omega
        =
        H_E^T\sigma H_\Omega,
        \qquad
        \sigma\ge_{\mathrm e}0.
        \label{eq:simplex-lp}
\end{equation}
The dimensions are
\begin{equation}
        I_E^T G I_\Omega\in\reals^{d_E\times d_\Omega},
        \qquad
        H_E^T\sigma H_\Omega\in\reals^{d_E\times d_\Omega}.
\end{equation}
Thus both sides represent the same bilinear pairing between accessible effects
and accessible states.

When Eq.~\eqref{eq:simplex-lp} is feasible, the Born statistics generated by all
pairs $(\rho,e)\in\OmegaSet\times\E$ admit a generalized-noncontextual,
equivalently simplex-embeddable, explanation.


When Eq.~\eqref{eq:simplex-lp} is feasible, the Born statistics generated by all pairs \((\rho,e)\in\OmegaSet\times\E\) admit a simplex-embeddable classical explanation for the specified fragment.  This is the sense in which the fragment is classical and ``noncontextual'' in the simplex test.  It should not be confused with the stronger KS requirement of a global two-valued valuation on an independently specified hypergraph.


If it is infeasible, one may
introduce an operational noise map
\begin{equation}
        D:\mathcal A\to\mathcal A
\end{equation}
on the state side and minimize $r$ subject to
\begin{equation}
        \begin{aligned}
        I_E^T G\bigl[(1-r)\operatorname{Id}_{\mathcal A}+rD\bigr]I_\Omega
        &= H_E^T\sigma H_\Omega,\\
        \sigma&\ge_{\mathrm e}0,\qquad 0\le r\le1.
        \end{aligned}
        \label{eq:noisy-lp}
\end{equation}
Here $\operatorname{Id}_{\mathcal A}$ denotes the identity map on the ambient
Hermitian-operator coordinate space.  The optimum $r$ is a robustness against
the specified noise.  In the calculations below, $D$ is always the completely
depolarizing state-side map
\begin{equation}
        \rho\longmapsto \tr(\rho)\frac{\id_{d_H}}{d_H},
        \label{eq:depolarizing-map}
\end{equation}
where $d_H=\dim\HH$.

Four qualifications fix the scope of this test.
First, the linear program used here is not a hypergraph-coloring or two-valued-state test.
In the restricted projector-cone version, its input is the set of projectors used as preparations and effects,
together with their Born pairing.
 The context structure is therefore seen only through the operator geometry of these projectors:
orthogonality, linear dependence, and relations such as \(\sum_{v\in C}P_v=\id\) when a complete context is present.
The linear program does not, however, impose the hypergraph rule that a deterministic valuation must select exactly one atom in each block.  Nor does it include, unless added explicitly, the full operational measurement structure consisting of unit effects, complements, coarse grainings, and context-induced operational equivalences.  The robustness values reported below should therefore be read as projector-cone robustness values, not as invariants of the abstract hypergraph alone.


Second, simplex embeddability is weaker than simpliciality.  A generalized
probabilistic theory or an accessible fragment may embed into a classical
simplex even though its own state space is not itself a simplex.  This
distinction is central to the generalized-probabilistic-theory analysis of
Schmid \emph{et al.}~\cite{schmid-2021-simplex}.  For accessible fragments,
cone equivalence is especially important: cone-equivalent fragments are either
both classically explainable or both not, because the embedding problem depends
on the generated cones~\cite{selby-2023-fragments}.

Third, Eq.~\eqref{eq:simplex-lp} is not the same object as the convex hull of
two-valued states of an orthogonality hypergraph.  The two-valued-state hull
appears as a sharp, deterministic shadow of the operational problem when the
fragment contains the eigenstate preparations, sharp projective measurements,
and operational equivalences that force deterministic responses on the relevant
ontic support.  This is the bridge to KS reasoning, not an identity
of the two frameworks.

Finally, the computations reported in Table~\ref{tab:results} use a deliberately
restricted \emph{projector-cone} version of the simplex test.  The listed
rank-one projectors are used as preparations and as effects, and the cone
generated by those projectors is tested against state-side depolarizing noise.
The unit effect, complementary coarse-grainings such as $\id-P$, and operational
equivalences induced by complete measurement contexts are not added unless they
are generated by the same projector cone.  Consequently, the block structure of
a hypergraph is not fully visible to this particular linear program: two
presentations with the same projector set but different groupings into contexts
define the same projector-cone problem.  This is not the full
Schmid--Selby--Wolfe--Kunjwal--Spekkens operational scenario; it is a controlled
fragment intended to compare the cone generated by several canonical ray sets.
The thresholds below should therefore not be compared unqualifiedly with
thresholds obtained from implementations that include the full set of
coarse-graining equivalences.

This is where the relative character of the taxonomy enters.  In the full
operational treatment, adding the unit effect, complements, mixtures,
context-induced operational equivalences, or local implementations of product
observables may turn a classically harmless restricted fragment into a
nonclassical one.  Conversely, forgetting such data may collapse a multi-context
or algebraic obstruction to a simplex-classical probability table.  The robust
claim is not that the three layers are permanently separated for every
enrichment of the data, but that they test different retained structures and
therefore need not agree before those structures are explicitly fixed.


\section{The simplicial-sandwich criterion and useful diagnostics}
\label{sec:sandwich}

The exact simplex criterion can be stated compactly as follows.  A finite
prepare-and-measure fragment is classical at the simplex level exactly when its
accessible Born pairing factors through a finite classical orthant.  In the
coordinate convention of Section~\ref{sec:simplex-test}, this Born pairing is
the bilinear form
\begin{equation}
        (y,x)\longmapsto y^T I_E^T G I_\Omega x.
\end{equation}
Thus, after quotienting by operational equivalences, there must exist a finite
ontic set $\Lambda$ and positive linear maps
\begin{equation}
        \mu:S_\Omega\longrightarrow \reals_+^\Lambda,
        \qquad
        \xi:S_E\longrightarrow \reals_+^\Lambda,
\end{equation}
such that, for all included preparations and effects,
\begin{equation}
        \tr(e\rho)
        =
        \sum_{\lambda\in\Lambda}\xi_e(\lambda)\mu_\rho(\lambda),
        \label{eq:orthant-factorization}
\end{equation}
with the usual normalization and response constraints.  Thus the classical
simplex appears as the normalized slice of the orthant $\reals_+^\Lambda$.
In this language, simplex nonclassicality is precisely the failure of the
positive factorization in Eq.~\eqref{eq:orthant-factorization} for the specified
fragment.

Equivalently, once the effect cone is identified with functionals on the
accessible state space through the Born pairing $I_E^T G I_\Omega$, the problem
may be pictured as a sandwiching problem.  One asks whether there is a
simplicial cone $K$ such that
\begin{equation}
        \Cone(\Omega_A)\subseteq K\subseteq \Cone(E_A)^*,
        \label{eq:simplicial-sandwich}
\end{equation}
where the dual cone is taken with respect to the same Born pairing, not with
respect to an arbitrary Euclidean dot product.  The cone $K$ is the unnormalized
classical state cone of a hidden-variable model, while $K^*$ supplies the
positive response functionals.  In the facet representation this sandwich
condition is exactly the nonnegative matrix equation \eqref{eq:simplex-lp};
with noise it becomes Eq.~\eqref{eq:noisy-lp}.  By Farkas duality, failure of
the sandwich is witnessed by a separating linear functional, equivalently by a
noncontextuality inequality for the chosen operational table.

This geometric picture also explains why the simplex test is not a graph invariant.  The orthogonality hypergraph records exclusivity and context incidence, but the simplex linear program sees the accessible cones, the Born bilinear form, and the operational equivalences.  Metric information such as $|\langle\psi|\phi\rangle|^2$ matters.  Conversely, as has been noted before, context grouping is invisible to a projector-only cone test unless it is encoded by additional effects, coarse grainings, or equivalence constraints.  The exact question is therefore always: can this specified operational fragment, with these equivalences, be sandwiched by a simplicial cone?

The following diagnostics are useful, but they should not be mistaken for necessary-and-sufficient graph criteria.  First, a fragment whose accessible state cone is already simplicial is normally classically harmless; non-simpliciality or overcompleteness of the generated cone is a warning sign.  But merely having many extreme rays is not by itself an obstruction, since a nonsimplicial cone may still lie inside a larger simplicial cone compatible with the effects.  Second, the effects must resolve the nonsimplicial geometry.  Coarse or incomplete effects can make a real obstruction invisible.  Third, the constraints must close a nontrivial consistency loop: cyclic exclusivity,  TIFS or KS implication chains, parity structures after spectral refinement, or equal-mixture preparation equivalences.  Without such a loop the hidden-variable sample space often has enough freedom to absorb the transition probabilities.

Three common mechanisms are worth separating.  An \emph{overlap or cyclic-exclusivity obstruction} appears when nonorthogonal quantum states and exclusive effects must be reconciled around a closed cycle; the completed pentagon is the canonical low-dimensional calibration.  A \emph{determinism obstruction} appears when eigenstate preparations and sharp effects force response functions to take values $0$ or $1$ on large parts of the ontic space; in sufficiently interlocked configurations this is the operational shadow of KS, TIFS, and related two-valued-state constraints.  A \emph{convex-decomposition obstruction} appears when different accessible mixtures represent the same operational preparation and hence must be assigned the same ontic distribution.  This last mechanism is the usual preparation-contextuality mechanism and can occur in small scenarios if the relevant mixture equivalences are included.

These mechanisms illuminate the examples below.  The firefly logic \(L_{12}\)~\cite{cohen} realizes complementarity but does not produce a simplex obstruction in the restricted projector-cone fragment~\cite{schaller-92,svozil-2018-b}.  The completed Klyachko--Can--Binicio\u{g}lu--Shumovsky pentagon already exhibits a cyclic probability-level gap, although it is not KS uncolorable~\cite{Klyachko-2008,wright:pent}.  Meager-state configurations expose deterministic two-valued-state defects in different ways~\cite{tkadlec-96}.  The Yu--Oh 13-ray and completed 25-ray configurations~\cite{Yu-2012}, the Cabello--Estebaranz--Garc\'ia-Alcaine 18--9 set~\cite{cabello-96}, and the Peres--Mermin square together with the 24--24 completion~\cite{peres111,mermin90b,peres-91,Pavicic-2005,svozil-2024-convert-pra-externalfigures} give increasingly constrained projector-cone fragments whose Born pairings fail simplex embeddability without noise.  GHZ is different: if one retains only the common joint spectral projectors of the four commuting global product observables, the fragment is a single Boolean context and hence simplex-classical.  Its nonclassicality enters through the additional requirement that values of global products factor into pre-existing local Pauli values, a product-preservation constraint not present in the projector-only sandwich~\cite{ghz,mermin90b,svozil-2024-convert-pra-externalfigures}.


\section{Operator-valued simplex tests and masked contexts}
\label{sec:operator-valued}

An ``operator-valued simplex'' is not a single unambiguous object.  Its meaning depends on what is taken as the operational data and what constraints are imposed on the deterministic classical vertices.

First, if an operator is used only as shorthand for its spectral projectors, then there is no essential change from the projector simplex.  A normal observable
\begin{equation}
        A=\sum_i a_iP_i
\end{equation}
contributes expectation values
\begin{equation}
        \langle A\rangle_\rho=\sum_i a_i\,\tr(\rho P_i),
\end{equation}
which are linear post-processings of the spectral-outcome probabilities.  If the projectors $P_i$ are already included as effects, adding $A$ itself changes only the coordinates of the data table.

Second, if one keeps only a degenerate operator and does not include enough commuting observables to resolve its eigenspaces, then the simplex test becomes coarser.  Some distinctions between projectors are hidden.  This can make simplex embeddability easier and may miss contextuality that appears after maximal refinement.  The matrix-pencil method addresses precisely this issue: a suitable linear combination of mutually commuting degenerate operators reveals the common spectral projectors and hence the underlying context~\cite{svozil-2024-convert-pra-externalfigures}.

Third, if the classical model is required to assign values to operators while preserving functional and product relations, the test changes substantially.  The deterministic vertices are no longer arbitrary assignments to spectral outcomes; they must also satisfy relations such as Eq.~\eqref{eq:product-rule}.  This is the setting of Peres--Mermin and GHZ parity arguments.  In such cases the obstruction may be algebraic even when a projector-only simplex for a single context would be trivial.

Consequently, an operator-valued simplex can mean at least three different things: a simplex for expectation values, a simplex for spectral-outcome probabilities, or a simplex whose vertices are operator value assignments constrained by functional calculus.  Only the last one captures GHZ-type product contradictions.  The first two are operational probability tests; the last is an operator-functional test.  This is the sense in which GHZ, although a single context after diagonalization, still violates classical expectations: the violation is not in the Boolean algebra of its joint spectral projectors but in the attempted decomposition of its global product observables into context-independent local elements of reality.


\subsection{Consecutive global measurements versus local factorizations}
\label{subsec:ghz-consecutive}

One might object that the four GHZ product observables commute.  Why not simply measure them consecutively on the same three-particle system?  In notation, the four observables are
\begin{align}
        A_1&=X_1X_2X_3,&
        A_2&=X_1Y_2Y_3,\nonumber\\
        A_3&=Y_1X_2Y_3,&
        A_4&=Y_1Y_2X_3 .
\end{align}
This consecutive protocol is possible in principle.  Since the $A_i$ commute, their spectral projections commute.  An ideal L\"uders measurement of one product observable collapses the state only to an eigenspace that is invariant under all the others.  Subsequent ideal measurements refine the state to a common eigenspace and do not erase the previously registered eigenvalues.  Equivalently, one may perform the joint projective measurement with atoms
\begin{equation}
        P_{\boldsymbol\varepsilon}
        =\prod_{i=1}^4 \frac{\id+\varepsilon_i A_i}{2},
        \qquad
        \varepsilon_i\in\{\pm1\},
        \qquad
        \varepsilon_1\varepsilon_2\varepsilon_3\varepsilon_4=-1,
        \label{eq:ghz-joint-pvm}
\end{equation}
where only eight sign patterns are nonzero because $A_1A_2A_3A_4=-\id$.  Thus a collective consecutive protocol registers four repeatable outcomes whose product is always $-1$.

\begin{remark}[Commuting observables and registered outcomes]\label{rem:commuting-lueders}
The preceding statement uses the standard finite-dimensional spectral fact that
commuting self-adjoint observables admit compatible L\"uders measurements.
Let $A$ and $B$ be self-adjoint operators with $[A,B]=0$, and let
$P_a$ and $Q_b$ be their spectral projectors.  Since spectral projectors
are obtained from the operators by functional calculus, $[P_a,Q_b]=0$ for
all eigenvalues $a,b$.  Hence $B P_a=P_aB$, so the subspace
$P_a\HH$ is invariant under $B$.  Equivalently,
\begin{equation}
        P_a\HH=\bigoplus_b P_aQ_b\HH,
        \label{eq:joint-eigenspace-refinement}
\end{equation}
up to zero summands.  Thus the collapse caused by registering the outcome
$a$ of $A$ does not cut across the eigenspace decomposition of $B$; it only
selects a direct sum of joint eigenspaces.

For a L\"uders measurement of $A$ followed by a measurement of $B$, the
sequential joint probability is
$$
        p(a\ \mathrm{then}\ b)
         =\tr\bigl(Q_b P_a\rho P_a\bigr)
         =\tr\bigl(P_aQ_b\rho\bigr),
        %\label{eq:sequential-equals-joint}
$$
which is precisely the probability assigned by the simultaneous joint projection-valued measure (PVM)
$\{P_aQ_b\}_{a,b}$.  A later measurement of $B$ may refine the post-measurement
state within $P_a\HH$, but it cannot invalidate the previously registered
value $a$.  This proof applies to the degenerate observable itself.  It
does not apply to an arbitrary finer apparatus that, while reporting the
same coarse outcome $a$, measures additional degrees of freedom inside
$P_a\HH$ which need not commute with $B$.
\end{remark}

This collective protocol does not, by itself, yield a GHZ contradiction.  It is just one joint measurement, hence one Boolean context.  A classical hidden-variable model for it is immediate: the ontic state is one of the eight allowed sign patterns $\boldsymbol\varepsilon=(\varepsilon_1,\varepsilon_2,\varepsilon_3,\varepsilon_4)$ satisfying $\prod_i\varepsilon_i=-1$, and the measurement reads off its four coordinates.  No product rule for local factors is used, and no inconsistency follows.

The GHZ contradiction appears only if the same global product values are required to arise from pre-existing local Pauli values,
$v(X_i),v(Y_i)\in\{\pm1\}$,
$v(X_1Y_2Y_3)=v(X_1)v(Y_2)v(Y_3)$,
with analogous equations for the other three products.  Multiplying the four classical equations gives $+1$, because every local factor occurs twice.  Quantum mechanically the product of the four commuting global operators is $-\id$.  The collective measurement of Eq.~\eqref{eq:ghz-joint-pvm} reveals the global products, but it never reveals the six local quantities $X_i,Y_i$ whose context-independent coexistence is being assumed.  In that precise sense the collective protocol is factor-blind.

The local GHZ protocol operationalizes a different refinement.  In each run the three separated parties choose one of the four compatible triples
        $\{X_1,X_2,X_3\}$,
        $\{X_1,Y_2,Y_3\}$,
        $\{Y_1,X_2,Y_3\}$,
        $\{Y_1,Y_2,X_3\}$,
and multiply the local outcomes.  These four triples are not jointly measurable by local sharp measurements, since on each particle $X_i$ and $Y_i$ anticommute.  Locality, or equivalently the Einstein--Podolsky--Rosen (EPR) counterfactual step in the usual GHZ reasoning, is what motivates assigning the same value to $X_i$ or $Y_i$ independently of which compatible triple is chosen.  The contradiction therefore concerns the attempted identification of two operationally different realizations of a product observable: a collective entangled measurement of the product as a whole, and a local measurement of its factors followed by multiplication.


There is a final technical caveat.  A degenerate observable does not determine a unique measurement apparatus.  A badly chosen implementation may disturb states inside a degenerate eigenspace by effectively measuring additional, possibly incompatible, degrees of freedom.  That is not a failure of commutativity of the abstract observables; it is a different physical refinement of the degenerate measurement.  For the ideal consecutive protocol considered here, one assumes L\"uders instruments or the common joint PVM, in which case the registered product outcomes are compatible and repeatable.  What remains nonclassical is not the sequence of global product outcomes, but the extra product-preserving identification of those outcomes with jointly pre-existing local values.


\subsection{Single-context immunity, unmasking, and the chromatic analogy}
\label{subsec:ghz-unmasking}

The preceding discussion also fixes a possible misconception about the scope of simplex tests.  A closed fragment consisting of one joint spectral measurement is always simplex embeddable.  Its ontic states may simply be the atoms of that Boolean algebra, and an arbitrary preparation is represented by the corresponding probability distribution over those atoms.  Thus the collective GHZ product fragment, taken by itself, has no projector-simplex obstruction.  In this limited and precise sense, the GHZ eigensystem is \emph{simplex immune}.

The qualifier is essential.  One should not conclude that operator-valued arguments can never enter a simplex analysis.  They can do so whenever the operator proof is unfolded into a multi-context projector fragment, or whenever the operational scenario includes the local factors and their compatibility relations.  Peres--Mermin supplies the clearest example: its algebraic parity proof, once the degenerate observables are resolved by simultaneous diagonalization or matrix pencils, yields a genuine multi-context projective configuration, and that configuration can be subjected to two-valued-state and simplex tests.  The GHZ collective context does not itself have this property; the local GHZ experiment has a different operational structure, namely four incompatible local settings together with the product-rule identification of shared local factors.

The operative distinction is therefore not whether an argument is operator-valued or simplex-theoretic in the abstract.  The right question is whether the retained fragment is merely a single collective spectral context or an unmasked multi-context fragment.  The former is classically representable as a Boolean block.  The latter may display ordinary projector contextuality, simplex nonembeddability, Bell--GHZ nonlocality, or an operator-functional parity contradiction, depending on which structures are retained in the fragment.

This also clarifies the relation to chromatic contextuality.  GHZ and strong chromatic contextuality are similar only in the negative sense that neither is exhausted by a probability-table simplex test.  Positively they are different obstructions.  Chromatic contextuality is combinatorial: it asks whether a fixed nondegenerate spectrum, or equivalently a strong coloring, can be assigned globally across a hypergraph.  GHZ is algebraic: it asks whether values can preserve products of commuting observables and their local factors.  State-independent parity proofs such as Peres--Mermin are closer in spirit to chromatic obstructions because they express a global spectral inconsistency independent of the input state, but the mechanism is still operator multiplication rather than hypergraph coloring.  Thus both belong outside the simplex-probability axis, but they lie on different non-simplex axes: one incidence/coloring-theoretic, the other operator-algebraic.

\section{Why two contexts need not, by themselves, reveal nonclassicality}
\label{sec:two-contexts}

Before turning to calibration examples, it is useful to remove one smaller temptation: the idea that the passage from one context to two already separates quantum vector probabilities from classical convex probabilities.  That intuition is correct at the level of the full theories but misleading for a small operational fragment.  A single orthonormal context is an ordinary probability simplex.  Two nonintertwining contexts are, operationally, a pair of simplexes connected by a stochastic transition matrix.

Let
\begin{equation}
        A=\{a_1,\ldots,a_d\},\qquad B=\{b_1,\ldots,b_d\}
\end{equation}
be two orthonormal bases.  The quantum transition matrix
\begin{equation}
        T_{ji}=|\langle b_j|a_i\rangle|^2
\end{equation}
is doubly stochastic.  For the restricted fragment consisting only of preparations in $A\cup B$ and measurements in $A\cup B$, this matrix can be reproduced by a classical random channel.  One may take ontic states to be pairs $\lambda=(i,j)$, where $i$ is the predetermined answer to context $A$ and $j$ is the predetermined answer to context $B$.  Preparing $a_i$ means loading the classical ensemble with the states $(i,1),\ldots,(i,d)$ in proportions $T_{1i},\ldots,T_{di}$.  A measurement of $A$ reads the first coordinate and therefore returns $i$ with certainty; a measurement of $B$ reads the second coordinate and returns the Born transition probabilities.  Preparations $b_j$ are treated symmetrically.

This construction is not a classical simulation of the full Hilbert-space theory.  It does not reproduce arbitrary third bases, interference phases, sequential state update, or the entire lattice of subspaces.  Nor is it a theorem that every possible two-context operational scenario is simplex embeddable: if additional preparations, operational equivalences, or linear dependences among preparations and effects are imposed, preparation contextuality may already be witnessed in very small scenarios, including scenarios with two binary measurements~\cite{pusey-2018}.  The limited claim here is that a bare transition table between two bases, with no further operational equivalences beyond normalization, can be represented as an ordinary stochastic channel.  Nonclassicality becomes visible when one asks for a single classical sample space to satisfy enough compatibility constraints around a nontrivial web of pasted contexts, or when additional equivalences constrain the preparations and effects.  The obstruction is not the mere fact of having two contexts, but the consistency requirements imposed on them.

This is precisely the lesson of partition logics.  A partition logic is obtained by pasting Boolean algebras arising from partitions of a finite set.  Its atoms can be represented as subsets of a common classical sample space, and probabilities are ordinary convex mixtures over dispersion-free states.  Equivalently, the same structures can be modeled by Wright-style generalized urns or by the initial-state identification problem for finite deterministic Moore or Mealy automata.  In a generalized urn, a ball type is an ontic state and different colors correspond to incompatible experimental contexts; looking through one color filter reveals one partition while hiding the others.  In the automaton model, the initial state is the ontic state and different input symbols induce different observable partitions of the state set.  These models realize complementarity without value indefiniteness.

The contrast with a Hilbert-space realization is then not in the exclusivity graph alone.  The same graph may support a set-theoretic partition representation, a faithful orthogonal vector representation, or even more exotic weights~\cite{svozil-2018-b}.  The probability type is determined by the chosen physical representation and the operational fragment, not by the bare pasting diagram alone.  Thus the statement ``complementarity does not imply contextuality'' is not merely philosophical: it is a concrete modeling fact.

\section{Calibration examples: firefly logic and the pentagon}
\label{sec:calibration}

The first calibration example is the smallest nontrivial pasting, the firefly logic $L_{12}$ is a 5--2 vertex--context set, obtained by pasting two three-atomic Boolean algebras along one common atom~\cite{cohen},
\begin{equation}
        \{a,b,c\},\qquad \{a,d,e\}.
\end{equation}
A faithful qutrit realization is
\begin{equation}
        a=(1,0,0),\quad b=(0,1,0),\quad c=(0,0,1),
\end{equation}
\begin{equation}
        d=(0,1,1),\quad e=(0,1,-1).
\end{equation}
The partition-logic representation is even more elementary.  On the finite set $S_5=\{1,\ldots,5\}$ one may use the two partitions
\begin{equation}
        \bigl\{\{2,3\},\{4,5\},\{1\}\bigr\},\qquad
        \bigl\{\{1\},\{3,5\},\{2,4\}\bigr\},
        \label{eq:l12-partition}
\end{equation}
which paste along the singleton atom $\{1\}$.  This is a generalized urn or finite-automaton realization of the same logic.  It is non-Boolean as a pasted structure, but it is still classically representable.  Its role here is diagnostic: it is the first place where contexts are not simply disjoint, yet simplex nonembeddability should not be expected merely from that fact.

The next calibration example is the pentagon, or pentagram logic.  In the Klyachko--Can--Binicio\u{g}lu--Shumovsky (KCBS) form one starts with five qutrit rays $v_i$ satisfying cyclic orthogonality
\begin{equation}
        v_i\perp v_{i+1}\qquad (i\;\mathrm{mod}\;5).
\end{equation}
A convenient realization is
\begin{equation}
        v_k=\left(
        \cos\frac{4\pi k}{5},\;
        \sin\frac{4\pi k}{5},\;
        \sqrt{\cos\frac{\pi}{5}}
        \right),\qquad k=0,\ldots,4.
        \label{eq:kcbs-rays}
\end{equation}
Completing each orthogonal pair by $w_i=v_i\times v_{i+1}$ gives five complete qutrit contexts
\begin{equation}
        \{v_i,w_i,v_{i+1}\},\qquad i=0,\ldots,4,
        \label{eq:pentagon-contexts}
\end{equation}
and hence a 10-ray pentagon (10--5 vertex--context set).

The pentagon sits between $L_{12}$ and the Specker bug discussed later.  It is still partition-logically representable: the five cyclically intertwined contexts support 11 two-valued states, from which a partition logic over $S_{11}=\{1,\ldots,11\}$ can be reconstructed.  Nevertheless, the pentagon already displays a difference between classical convex probabilities and Born/Lov\'asz-type vector probabilities.  On the five outer atoms, the classical convex hull yields the familiar pentagon bound $\sum_i p_i\le2$, whereas the optimal faithful orthogonal vector representation gives the quantum value $\sqrt5$.  This is the first useful calibration point at which cyclic exclusivity exposes a quantitative difference, though still not a KS impossibility.

\section{Combinatorial contextuality: two-valued states and meagerness}

Let $H=(V,\C)$ be a finite $d$-uniform orthogonality hypergraph.  The vertices $v\in V$ represent atoms, usually rank-one projectors, and each context $C\in\C$ is a $d$-element set of mutually orthogonal atoms corresponding to a maximal sharp measurement.  A faithful orthogonal representation in $\complex^d$ or $\reals^d$ assigns to every vertex a ray $|v\rangle$ such that vertices are orthogonal precisely when the hypergraph says they are, and every context resolves the identity:
\begin{equation}
        \sum_{v\in C} \Pi_v = \id,
        \qquad \Pi_v=\frac{|v\rangle\langle v|}{\langle v|v\rangle}.
\end{equation}

\begin{definition}[Two-valued state]
A two-valued state on $H$ is a map $s:V\to\{0,1\}$ such that
\begin{equation}
        \sum_{v\in C}s(v)=1
        \qquad \text{for every } C\in\C.
        \label{eq:tvs}
\end{equation}
We denote the set of all such states by $\TVS(H)$.
\end{definition}

The deterministic valuation polytope is
\begin{equation}
        \Ptwo(H)=\conv\TVS(H)\subseteq[0,1]^V.
\end{equation}
It is useful to distinguish several ways in which $\TVS(H)$ can be deficient.  Tkadlec called attention to hypergraphs whose state spaces are empty, not unital, not separating, or not full~\cite{tkadlec-96}; these notions are standard in the quantum-logic literature~\cite{pulmannova-91,schaller-92,svozil-tkadlec}.

\begin{definition}[Unital, separating, full]
Let $\TVS(H)$ be the set of two-valued states on $H$.
It is called
\begin{enumerate}
\item \emph{unital} if for every vertex $a\in V$ there exists $s\in\TVS(H)$ with $s(a)=1$;
\item \emph{separating} if for every distinct pair $a,b\in V$ there exists $s\in\TVS(H)$ with $s(a)\ne s(b)$;
\item \emph{full} if for every nonorthogonal pair $a,b\in V$ there exists $s\in\TVS(H)$ with $s(a)=s(b)=1$.
\end{enumerate}
In the usual orthomodular-poset formulation, where atoms sit inside a poset with zero, unit, and orthocomplementation, fullness implies separation and separation implies unitality.  In the atom-only hypergraph formulation used here these implications require the corresponding reconstruction assumptions.  We therefore treat unitality, separation, and fullness as separate diagnostics and do not rely on the hierarchy without those assumptions.
\end{definition}

\begin{definition}[Meager two-valued-state space]
A hypergraph has a meager two-valued-state space if $\TVS(H)$ is empty, nonunital, nonseparating, or nonfull.  This is an umbrella term; the four cases have different operational meanings.
\end{definition}

Each failure mode corresponds to a different linear constraint on $\Ptwo(H)$.  If $\TVS(H)=\emptyset$, then $\Ptwo(H)=\emptyset$.  This is the KS case.  If $\TVS(H)$ is nonunital at $a$, then
\begin{equation}
        x_a=0 \qquad \text{for all }x\in\Ptwo(H).
        \label{eq:nonunital-face}
\end{equation}
If it is nonseparating for $a\ne b$, meaning $s(a)=s(b)$ for all $s\in\TVS(H)$, then
\begin{equation}
        x_a=x_b \qquad \text{for all }x\in\Ptwo(H).
        \label{eq:nonsep-face}
\end{equation}
If it is nonfull for a nonorthogonal pair $a,b$, meaning no $s$ has $s(a)=s(b)=1$, then
\begin{equation}
        x_a+x_b\le1 \qquad \text{for all }x\in\Ptwo(H).
        \label{eq:nonfull-face}
\end{equation}
A true-implies-false (TIFS) structure, such as the Specker bug, is the stronger directed version: for a nonorthogonal pair $a,b$,
\begin{equation}
        s(a)=1\Longrightarrow s(b)=0,
        \label{eq:tifs}
\end{equation}
for all admissible $s$.  Equation~\eqref{eq:tifs} implies that the face $x_a=1$ of the two-valued-state polytope is contained in $x_b=0$.

\section{How meagerness becomes operational}

The preceding notions are combinatorial: they concern the possible two-valued
states on a hypergraph.  To become operational, however, such defects must be
visible in the actual preparations and effects included in the fragment.  A
logical obstruction that is present in the hypergraph may remain invisible if
the fragment is too poor to test it.

Let $H$ have a faithful orthogonal representation, with vertex projectors
$\Pi_v$.  For a preparation $\rho$, write
\begin{equation}
        p^\rho_v=\tr(\rho\Pi_v)
\end{equation}
for the Born probability assigned to the vertex effect $\Pi_v$.  The
deterministic two-valued-state polytope
\begin{equation}
        \Ptwo(H)=\conv\TVS(H)
\end{equation}
contains exactly the probability assignments obtainable as classical mixtures
of admissible two-valued states.  Thus, if an accessible Born vector
$p^\rho$ lies outside $\Ptwo(H)$, then no deterministic two-valued-state
model can reproduce that part of the fragment.

The different forms of meagerness become visible in different ways.
If the two-valued states are nonunital at an atom $a$, then every admissible
two-valued state assigns
\begin{equation}
        s(a)=0.
\end{equation}
Consequently every classical mixture in $\Ptwo(H)$ also assigns probability
zero to $a$.  But if the fragment allows the eigenstate preparation
$\rho=\Pi_a$ and the effect $\Pi_a$, quantum theory gives
\begin{equation}
        \tr(\Pi_a\Pi_a)=1.
\end{equation}
The atom that is never true in the two-valued-state model is obtained with
certainty in the corresponding quantum eigenstate.  Nonunitality is then
directly operationally visible.

If the two-valued states are nonseparating for two distinct atoms $a$ and
$b$, then every admissible valuation identifies them:
\begin{equation}
        s(a)=s(b).
\end{equation}
Every classical mixture therefore satisfies $x_a=x_b$.  This becomes
operationally relevant only if the Born statistics can distinguish the two
effects, that is, if the fragment contains some preparation $\rho$ such that
\begin{equation}
        \tr(\rho\Pi_a)\ne \tr(\rho\Pi_b).
\end{equation}
In that case the classical valuation model merges two outcomes that the quantum
statistics separate.  If no such preparation is included, the nonseparability
remains present combinatorially but is not witnessed by the fragment.

Finally, a true-implies-false (TIFS) or nonfullness constraint becomes visible
through nonorthogonal overlap.  Suppose $a$ and $b$ are nonorthogonal, but
the two-valued states force
\begin{equation}
        s(a)=1\Longrightarrow s(b)=0.
\end{equation}
Classically, conditioning on $a$ being true then forces $b$ to be false.
Quantum mechanically, however, the eigenstate $\Pi_a$ gives
\begin{equation}
        \tr(\Pi_a\Pi_b)=|\langle a|b\rangle|^2>0
\end{equation}
whenever $a$ and $b$ are nonorthogonal.  Thus, if the fragment contains the
preparation $\Pi_a$ and the effects $\Pi_a,\Pi_b$, the TIFS implication is
operationally contradicted by the nonzero transition probability from $a$ to
$b$.

The important qualification is that these implications are fragment-dependent.
A meager hypergraph does not automatically fail every simplex test.  If a
nonunital atom is never prepared or measured, if a nonseparated pair is never
statistically distinguished, or if the endpoints of a TIFS gadget are absent
from the preparation/effect list, then the corresponding combinatorial defect
need not appear in the operational probability table.  Conversely, once the
relevant eigenstate preparations and effects are included, the valuation defect
becomes an operational obstruction, and the noisy simplex linear program
quantifies how much depolarizing noise is needed to hide it.

\section{Relation to generalized noncontextuality}

The deterministic two-valued-state analysis is closest to the traditional KS setting.  Generalized noncontextuality is broader: it does not assume determinism from the outset.  Nevertheless, determinism for sharp projective measurements is recovered on the support of eigenstate preparations when the operational equivalences include the relevant projective measurements and mixture decompositions.  This is the familiar route from operational noncontextuality to KS-style constraints, developed in modern form by Kunjwal and Spekkens~\cite{kunjwal-spekkens-2015}.

In this situation the two-valued-state polytope $\Ptwo(H)$ is the deterministic shadow of the operational simplex model.  Meagerness of $\TVS(H)$ supplies faces or equalities that the Born vectors may violate.  Depolarizing noise moves the Born vectors toward a more symmetric point, and the noisy simplex linear program computes the amount of noise needed to make the entire fragment simplex-embeddable.

The crucial qualification is that the simplex linear program is not merely ``the fractional coloring polytope.''  Its variables are the entries of $\sigma$ in Eq.~\eqref{eq:simplex-lp}, indexed by facets of the state and effect cones.  The number of such facets need not equal the number of vertices or contexts of the hypergraph.  In the Yu--Oh 13-ray computation below, there are 13 projectors but 24 facets of each cone, hence $\sigma$ is $24\times 24$.

\section{Combinatorial contextuality: chromatic completeness}

Chromatic contextuality imposes a different global requirement.  Let $H=(V,\C)$ be $d$-uniform.  A proper $d$-coloring is a map
\begin{equation}
        c:V\to\{1,\ldots,d\}
\end{equation}
such that every context contains every color exactly once.  Equivalently, for every color $k$ define
\begin{equation}
        s_k(v)=\begin{cases}
        1, & c(v)=k,\\
        0, & c(v)\ne k.
        \end{cases}
\end{equation}
Then each $s_k$ is a two-valued state and
\begin{equation}
        \sum_{k=1}^d s_k(v)=1
        \qquad \text{for every }v\in V.
        \label{eq:color-decomp}
\end{equation}
Thus chromatic colorability asks not merely for two-valued states, but for a complete decomposition of the unit assignment into $d$ admissible two-valued states.  The obstruction $\chi(H)>d$ means that no such global spectral-label decomposition exists.

This explains why chromatic contextuality is independent of ordinary valuation scarcity.  A hypergraph may possess a separating or even full set of two-valued states and nevertheless fail Eq.~\eqref{eq:color-decomp}.  Conversely, a KS hypergraph with no two-valued states is of course chromatically contextual, but the chromatic description is then not the most discriminating one: the stronger valuation obstruction is already present.



\section{Computational protocol}

The computational analysis fixes one convention for all examples.  The listed rays define rank-one projectors, and these projectors are used both as preparations and as effects.  Depolarizing noise is applied only on the state side.  The named hypergraphs determine the projector sets under consideration, and their block structure is used in the accompanying combinatorial analysis of two-valued states, colorability, and KS-type constraints.  In the projector-cone linear programs, however, the blocks are not imposed as additional valuation or measurement constraints.  In particular, the program does not separately enforce rules such as ``exactly one atom in each block,'' nor does it add the full family of unit effects, complements, coarse grainings, and context-induced operational equivalences associated with complete measurements.  The blocks enter only through the operator geometry of the projectors included in the cones, for example through orthogonality and linear relations among the projectors.  This convention is intentionally narrower than the full operational simplex-embeddability framework and is used here only to compare the cones generated by several canonical projector sets.


The computations use the non-orthonormal coordinate convention described in
Section~\ref{sec:simplex-test}.  For exact rational audits the program represents
real symmetric projectors in the raw coordinate basis
\begin{equation}
        (\rho_{11},\ldots,\rho_{dd},
        \rho_{12},\rho_{13},\ldots,\rho_{d-1,d}),
        \label{eq:raw-basis}
\end{equation}
rather than in a trace-orthonormal basis.  In this basis the trace pairing is
represented by the diagonal matrix
\begin{equation}
        G=\operatorname{diag}(1,\ldots,1,2,\ldots,2),
        \label{eq:raw-trace-metric}
\end{equation}
because off-diagonal real symmetric entries contribute twice to the trace
inner product.  Thus, if $e$ and $\rho$ denote raw coordinate vectors, then
\begin{equation}
        \tr(e\rho)=e^T G\rho.
        \label{eq:raw-trace-pairing}
\end{equation}
Equivalently, the code may store the dual effect vector $Ge$, whose ordinary
dot product with $\rho$ gives the same Born probability.  This convention avoids
square-root factors for integer rays and makes the rational cases suitable for
exact polyhedral conversion with \texttt{pycddlib}/cddlib over $\mathbb Q$.

The audit uses the accessible-fragment formalism.  Starting from full raw
coordinate matrices $V^F_\Omega$ and $V^F_E$, Stage~1 replaces them by
accessible coordinate matrices $V^A_\Omega$ and $V^A_E$ of full row rank,
together with inclusion matrices $I_\Omega$ and $I_E$.  These inclusion matrices
embed the accessible coordinates back into the full raw coordinate space, so
that
\begin{equation}
        V^F_\Omega=I_\Omega V^A_\Omega,
        \qquad
        V^F_E=I_E V^A_E
\end{equation}
up to the chosen basis convention.  Stage~2 computes facet descriptions
$H_\Omega$ and $H_E$ of the accessible cones.  Stage~3 solves the metric-aware
noisy simplex equation
\begin{equation}
        I_E^T G\bigl[(1-r)\operatorname{Id}_{\mathcal A}+rD\bigr]I_\Omega
        =
        H_E^T\sigma H_\Omega,
        \,
        \sigma\ge_{\mathrm e}0,
        \,
        0\le r\le1.
        \label{eq:stage3-audit}
\end{equation}
This is the same equation as Eq.~\eqref{eq:noisy-lp}, written after the
accessible reduction and in the raw-coordinate metric convention.

Candidate optima are found numerically, but the reported exact-primal rows are
then checked symbolically with rational $G,H_\Omega,H_E,r$, and $\sigma$.

The certificate status in Table~\ref{tab:results} distinguishes three cases.  ``Exact optimum'' means that both an exact primal certificate and an exact dual certificate have been verified.  ``Exact primal'' means that the displayed rational $r$ is an exact feasible value for Eq.~\eqref{eq:stage3-audit}; it is therefore an exact upper bound matching the numerical optimum, but exact optimality is not yet certified because a rational dual witness has not been recovered.  ``Numerical'' means that only a floating-point LP solution and residual check are available.  The completed pentagon is algebraic rather than rational in the chosen coordinates and is therefore presently treated by the numerical branch.  The Specker bug, a truncated form of the KS $\Gamma_1$, and $\Gamma_3$ of KS~\cite{kochen1} are also retained as numerical rows from the earlier high-precision computation.

\section{Examples and results}


\subsection{Calibration computations: \texorpdfstring{$L_{12}$}{L12} and the completed pentagon}

The firefly logic $L_{12}$ is included as a calibration for complementarity without simplex obstruction.  Its five projectors span a four-dimensional accessible state space inside the six-dimensional real symmetric qutrit trace space.  After the Stage~1 accessible-fragment reduction the exact rational audit finds five facets for each accessible cone and returns
\begin{equation}
        r_{L_{12}}=0,
\end{equation}
with both exact primal and exact dual certificates.  Thus $L_{12}$ is not merely a lower-rank exception to the full-rank trace-space code; it is an exactly simplex-embeddable accessible fragment, as expected from its partition-logic realization.

The completed pentagon contains 10 projectors and spans the full six-dimensional real symmetric qutrit trace space.  The computation verifies adjacent orthogonality, obtains five contexts of the form Eq.~\eqref{eq:pentagon-contexts}, finds 12 cone facets, and returns
\begin{equation}
        r_{C_5}\simeq0.119140568134.
        \label{eq:pentagon-result}
\end{equation}
The value is numerical only in the present audit because the coordinates involve algebraic quantities outside the rational cddlib backend.  It is small compared with the later KS-type examples, as expected: the pentagon is cyclic and already inequality-contextual in the KCBS sense, but it is still far weaker than a KS or nonunital obstruction.

\subsection{Some meager-state configurations}

The 13-ray Specker bug~\cite{2018-minimalYIYS} is TIFS, a true-implies-false gadget, and hence a nonfull two-valued-state example.
The KS $\Gamma_3$ configuration, with a nonseparating set of two-valued states.
We use a coordinatization of Tkadlec, who also exposes a the nonunital Schutte-type configuration~\cite{tkadlec-96}.  These examples supply the three nonempty forms of meagerness: nonfull, nonseparating, and nonunital.  In the projector-only simplex test, these particular ray sets yield different robustness values.  This shows operational inequivalence of the chosen fragments, not a universal monotone ordering of the three forms of meagerness.

For the Specker bug, the program uses 13 rays, obtains 26 facets, and returns
\begin{equation}
        r_{\mathrm{Specker\ bug}}\simeq0.31645570.
\end{equation}
For $\Gamma_3$, the input contains 26 listed rays but 25 unique projectors after identifying antipodal rays.  The cone has 230 facets and the LP returns
\begin{equation}
        r_{\Gamma_3}\simeq0.50219773.
\end{equation}
For Fig.~2, the 37-ray nonunital example has 276 facets.  The numerical optimum rationalizes to
\begin{equation}
        r_{\mathrm{nonunital}}=\frac{35}{64},
\end{equation}
and the rational audit verifies an exact primal certificate at this value.  A rational dual certificate has not yet been recovered, so the value is reported below as an exact feasible upper bound matching the numerical optimum, not as a formally certified optimum.  The three Tkadlec rows should not be read as universal invariants of the labels ``nonfull'', ``nonseparating'', and ``nonunital''.

\subsection{Yu--Oh: 13 rays versus 25-ray orthogonal completion}

The original Yu--Oh construction consists of the 13 qutrit rays
\begin{align}
&(1,0,0),(0,1,0),(0,0,1),\nonumber\\
&(0,1,1),(0,1,-1),(1,0,1),(1,0,-1),(1,1,0),(1,-1,0),\nonumber\\
&(1,1,1),(1,1,-1),(1,-1,1),(-1,1,1).
\end{align}
The projector fragment has 24 cone facets.  The simplex linear program gives
\begin{equation}
        r_{\YO13}=\frac38.
        \label{eq:yo13-result}
\end{equation}
In the facet ordering returned by the exact audit program, an optimal certificate can be chosen diagonal, with 18 nonzero diagonal entries and exact residual zero.

The 25-ray version used here is the orthogonal completion of the 13 rays: for every orthogonal pair among the original 13 rays, the cross-product ray completing the triad is added, and duplicates are removed projectively.  The completion contains 25 projectors, its cone has 96 facets, and the LP returns
\begin{equation}
        r_{\YO25}=\frac{21}{44}\simeq0.47727273.
        \label{eq:yo25-result}
\end{equation}
This increase over $3/8$ is expected: the enlarged projector set generates a different operational cone and supplies more Born constraints to be classically simulated.  The most recent rational audit verifies an exact primal certificate at $r=21/44$, matching the numerical optimum.  Since no exact dual witness has yet been recovered, the table marks this as exact-primal rather than exact-optimal.

\subsection{Cabello 18--9 and the Peres--Mermin 24--24 completion}

The Cabello--Estebaranz--Garcia-Alcaine 18--9 set is a four-dimensional KS configuration with 18 rays in 9 tetrads~\cite{cabello-96}.  The projector-only fragment is represented in the 10-dimensional real symmetric trace basis.  The cone generated by the 18 projectors has 146 facets.  The numerical optimum rationalizes to $1/3$, and the exact rational audit verifies an exact primal certificate,
\begin{equation}
        r_{18/9}=\frac13.
        \label{eq:cabello-result}
\end{equation}
As for Yu--Oh 25, exact optimality is not yet certified because a rational dual witness is still missing.

The Peres--Mermin square, analyzed through matrix pencils, yields the 24-24 completion.  The 24 vectors include the Cabello 18-ray set as a subset; the remaining rays supply the orthogonal completion associated with the full 24--24 Peres configuration~\cite{svozil-2024-convert-pra-externalfigures}.  In the projector-only simplex test, the 24 projectors generate a cone with 120 facets.  The numerical optimum rationalizes to $4/9$, and the exact rational audit verifies an exact primal certificate,
\begin{equation}
        r_{\mathrm{PM24}}=\frac49.
        \label{eq:pm24-result}
\end{equation}
Again, the increase is consistent with adding projectors and hence strengthening the operational fragment.  Exact dual certification of optimality remains open in the present audit.



\begin{table*}[t]
\caption{Depolarizing robustness for the restricted projector-cone prepare-and-measure fragments used in this paper.  The listed rays are used both as rank-one preparations and as rank-one effects.  The noise is the state-side depolarizing map $\rho\mapsto \tr(\rho)\id_d/d$.  ``Exact optimum'' means exact primal and exact dual certificates.  ``Exact primal'' means an exact feasible rational certificate matching the numerical optimum, but with no exact dual witness yet recovered.  ``Numerical'' means floating-point certificate only.}
\label{tab:results}
\resizebox{\textwidth}{!}{%
\begin{tabular}{lcccccc}
\toprule
Configuration & $d$ & projectors & rank & facets & $r$ & certificate status\\
\midrule
Firefly logic $L_{12}$ & 3 & 5 & 4 & 5 & $0$ & exact optimum\\
Completed pentagon $C_5$ & 3 & 10 & 6 & 12 & $0.119140568134$ & numerical/algebraic branch\\
Specker bug/TIFS, nonfull & 3 & 13 & 6 & 26 & $0.31645570$ & numerical\\
$\Gamma_3$, nonseparating & 3 & 25 & 6 & 230 & $0.50219773$ & numerical\\
Tkadlec Fig.~2, nonunital & 3 & 37 & 6 & 276 & $35/64$ & exact primal; dual open\\
Yu--Oh 13-ray projector fragment & 3 & 13 & 6 & 24 & $3/8$ & exact optimum\\
Yu--Oh 25-ray orthogonal completion & 3 & 25 & 6 & 96 & $21/44$ & exact primal; dual open\\
Cabello--Estebaranz--Garcia-Alcaine 18--9 & 4 & 18 & 10 & 146 & $1/3$ & exact primal; dual open\\
Peres--Mermin 24--24 completion & 4 & 24 & 10 & 120 & $4/9$ & exact primal; dual open\\
\bottomrule
\end{tabular}%
}
\end{table*}



\section{Interpretation of the table}

With these conventions fixed, the numbers in Table~\ref{tab:results} should be read narrowly.  They are not chromatic numbers and are not counts of two-valued states.  They are operational noise thresholds for the specific projector fragments described above.  They answer the following question: how much state-side depolarizing noise is needed before the Born bilinear form on the cones generated by these projectors admits a simplex embedding?

The calibration examples show why the number of contexts alone is not the issue.  The firefly logic is now included after the accessible-fragment reduction and has $r=0$ with exact primal and dual certificates; conceptually it is a classically representable partition logic.  The completed pentagon gives a small but positive numerical robustness, in agreement with its role as the first cyclic exclusivity example.

The Tkadlec examples show that three specific meager two-valued-state spaces give three distinct projector-cone robustness values.  The Specker-bug/TIFS obstruction is nonfull: it forbids certain nonorthogonal atoms from being jointly true in any admissible valuation.  The $\Gamma_3$ obstruction is nonseparating: the two-valued states fail to distinguish certain distinct propositions.  The nonunital example has an atom that is never assigned value one.  The numerical ordering in Table~\ref{tab:results} should not be read as a theorem that nonfullness is always weaker than nonseparation or nonunitality.  The configurations differ in size and geometry.  What the table establishes is only that these particular fragments are operationally inequivalent under the specified projector-cone LP.

The Yu--Oh comparison shows that adding projectors matters.  The 13-ray and 25-ray versions are related combinatorially, but they define different operational cones.  The simplex linear program detects this difference and returns different thresholds.  The 13-ray value is an exact optimum, while the 25-ray value has an exact primal certificate but no exact dual certificate yet.

The Cabello/Peres--Mermin comparison gives the analogous four-dimensional case.  The 18--9 set is a critical KS subset of the Peres 24-ray completion.  The full 24-ray completion contains the 18-ray set and imposes additional projector-cone constraints, increasing the threshold in this restricted test.  In both rows the displayed rational value is supported by an exact primal certificate matching the numerical optimum, but exact dual optimality remains open.

These examples illustrate the central message: logical meagerness, chromatic noncolorability, simplex nonembeddability, and operator-functional product obstruction are related but none is a mere restatement of the others.  The simplex linear program provides an operationally parametrized robustness for the selected probability fragment.  Quantum logic identifies the structural kind of valuation failure.  Chromatic contextuality identifies a global spectral-labeling obstruction.  Operator-valued parity arguments identify algebraic obstructions to product-preserving value assignments, and these may be invisible if one looks only at a single joint spectral context.


\section{Discussion}
\label{sec:discussion}

The preceding examples support the layered terminology proposed in the introduction.  \emph{Combinatorial contextuality} asks about global assignments on a pasted context structure.  Its invariants include the existence of two-valued states, unitality, separability, fullness, partition-logic embeddability, and strong colorability.  \emph{Simplex contextuality} asks whether a specified operational probability table embeds into a classical simplex.  Its invariant is not a coloring or a set of valuations but the feasibility, or noisy feasibility, of Eq.~\eqref{eq:noisy-lp}.  \emph{Operator-functional contextuality} adds a further algebraic requirement: deterministic values assigned to observables must respect functional and product relations among commuting operators.

The first two notions are not disjoint in the sense of mutually exclusive classes.  They often coexist.  KS configurations such as Cabello 18--9 and the Peres--Mermin/Peres 24-ray completion are nonclassical at the valuation level and simplex contextual for their natural projector fragments.  The point is rather that they are different predicates on different retained data.  If the full operational-equivalence structure is retained, some distinctions made by the restricted projector-cone test may disappear; if that structure is forgotten, a genuine operational obstruction may become invisible.  Thus the comparisons in Table~\ref{tab:results} should be read as projector-cone comparisons, while the taxonomy itself concerns the broader choice between incidence data, convex-operational data, and operator-algebraic data.

The operator-functional layer explains why operator-valued arguments cannot be identified with either of the first two layers.  If operators are replaced by their spectral projectors, the argument may reduce to a projective quantum-logic structure, as in the Peres--Mermin transcription to the 24-ray completion.  If the operators remain coarse-grained, the probability data may be weaker than the maximal projective refinement.  If product rules are imposed on deterministic operator values, new parity contradictions appear.  GHZ is the decisive example: the common eigensystem is a single Boolean context and is simplex-classical as such, but the product of the four global observables is $-\id$, whereas any assignment of pre-existing local factor values predicts $+1$.  The lesson is not that operator-valued arguments are forever invisible to simplex methods; it is that simplex visibility depends on whether the operator proof has been unpacked into a multi-context operational fragment.

Partition logics make the separation especially transparent.  They are classically representable by construction: generalized urns and deterministic finite automata give explicit classical models.  Yet some of the same pasting diagrams also have faithful orthogonal representations that support Born probabilities different from the convex hull of the dispersion-free states.  The pentagon is the canonical small example: its classical convex hull gives the pentagon bound, while a faithful vector representation gives the Lov\'asz/Born value.  Thus complementarity and non-Booleanity do not force contextuality; cyclic intertwinement is the first place where the probability calculus becomes visibly nontrivial.

The computational table should therefore not be read as a ranking of one scalar called contextuality.  The Specker bug, $\Gamma_3$, and Tkadlec's nonunital configuration realize different forms of meagerness.  Yu--Oh 13 and Yu--Oh 25 are related but generate different operational cones.  Cabello 18--9 and Peres--Mermin/Peres 24 show how adding the orthogonal completion can increase simplex robustness.  The LP thresholds quantify operational nonclassicality for the selected projector fragments; they do not replace either the combinatorial classification or the operator-functional parity analysis.

Chromatic contextuality fits naturally into the combinatorial side.  It is not merely the absence of two-valued states.  A strong $d$-coloring is a much more rigid object: it is a global decomposition of the unit assignment into $d$ valuation layers, corresponding to a globally reusable nondegenerate spectrum.  As the completed Yu--Oh and $G_{32}$ comparison shows, chromatic obstruction is also distinct from geometric coordinatizability.  A hypergraph may fail chromatic completeness while still having a faithful orthogonal representation; another may fail coordinatizability for purely projective-geometric reasons.

This leaves room for probability models beyond both classical convex hulls and quantum Born tables.  Additive, context-independent exotic weights satisfy the local Kolmogorov constraints on each context but need not be convex mixtures of dispersion-free states or Hilbert-space probabilities.  They occupy the probability layer, not the valuation/coloring layer.  For this reason, the safest organizing principle is not a single ladder from classical to quantum, but a grid: one axis records the combinatorial assignment structure, another records the probability model placed on it, and a third records algebraic constraints among operator-valued observables.


\section{Conclusion}

Combinatorial, simplex, and operator-functional contextuality are best understood as classicality tests applied at different levels.  Combinatorial contextuality asks whether a pasted context structure admits global valuations, separating partition representations, or strong colorings.  Simplex contextuality asks whether a specified operational probability table factors through a classical simplex, possibly after adding noise.  Operator-functional contextuality asks whether values assigned to observables can preserve the functional and product relations that the corresponding quantum operators satisfy.

In the simplex case, the most useful geometric picture is the failure of a simplicial sandwich: no classical simplicial cone can be inserted between the accessible preparation cone and the dual accessible effect cone while preserving the Born pairing and the declared operational equivalences.  This is why simplex contextuality is sensitive to the chosen fragment, to metric Born probabilities, and to coarse-graining or preparation-equivalence data, rather than merely to the abstract incidence hypergraph.

The results reported here show why all three notions are needed.  The firefly logic $L_{12}$ and partition-logical examples demonstrate complementarity without simplex obstruction.  The completed pentagon demonstrates the first cyclic probability-level separation between the classical convex hull and vector probabilities.  Tkadlec's examples distinguish nonfull, nonseparating, and nonunital meagerness.  Yu--Oh, Cabello 18--9, and Peres--Mermin/Peres 24 show how natural quantum projector fragments become simplex contextual with positive depolarizing thresholds.  GHZ-type arguments show a different phenomenon: a single common spectral context may be classically representable as a Boolean block, while the operator products whose eigensystem defines that block still contradict any product-preserving classical value assignment.  The comparison with chromatic contextuality is therefore only partial.  Both GHZ and chromatic obstructions can escape a probability-only simplex test, but GHZ does so through operator multiplication, whereas chromatic contextuality does so through a global incidence/coloring constraint.

The terminology therefore matters.  \emph{Combinatorial contextuality} covers two-valued-state scarcity, Boolean-embedding failure, and chromatic spectral-labeling obstruction.  \emph{Simplex contextuality} covers the operational failure of classical simplex embeddability.  \emph{Operator-functional contextuality} covers parity and product-rule contradictions among observables.  They coincide in some canonical quantum examples, but none reduces to the other.  Future work should make the numerical certificates reported here exact where possible, compare the simplex thresholds with the full family of classical, quantum, and exotic probability bodies associated with the same pasting diagrams, and formulate operator-valued simplex tests that explicitly distinguish expectation-value, spectral-outcome, and product-preserving versions.


\begin{acknowledgments}
This research was funded in whole or in part by the \textit{Austrian Science Fund (FWF)} [Grant \href{https://doi.org/10.55776/PIN5424624}{digital object identifier (DOI): 10.55776/PIN5424624}].
The author acknowledges TU Wien Bibliothek for financial support through its Open Access Funding Programme.

This text was partially created and revised with assistance from large language models. All content, ideas, and prompts were provided by the author.
\end{acknowledgments}


\appendix

\section{Computational audit data}
\label{app:audit}

This appendix records the computational status behind Table~\ref{tab:results}.  It is included because the distinction between numerical optima, exact primal certificates, and exact dual certificates is essential for interpreting the table.

The exact audit uses the accessible-fragment form of Eq.~\eqref{eq:stage3-audit}.  Integer or rational rays are converted to raw real-symmetric projector coordinates, effects are multiplied by the trace metric $G=\operatorname{diag}(1,\ldots,1,2,\ldots,2)$, and cddlib is used over $\mathbb Q$ to obtain rational facet matrices whenever possible.  A numerical linear program first proposes $r$ and $\sigma$.  The proposed rational values are then checked exactly by verifying
\begin{equation}
        H_E^T\sigma H_\Omega
        =I_E^T\bigl[(1-r)I+rD\bigr]I_\Omega,
        \qquad \sigma\ge_{\mathrm e}0.
        \label{eq:exact-primal-check}
\end{equation}
This is a primal feasibility certificate.  Exact optimality additionally requires a rational dual witness.  The current audit verified such dual witnesses only for $L_{12}$ and Yu--Oh 13.

The final Python audit returned the status summarized in Table~\ref{tab:appendix-audit}.  The table is intentionally diagnostic rather than a second robustness table: its purpose is to show which entries are exact optima, which have only exact primal feasibility certificates, and which remain numerical.  In the rows marked ``exact primal'', the rational value displayed in Table~\ref{tab:results} is certified as feasible by Eq.~\eqref{eq:exact-primal-check}; exact optimality would additionally require a rational dual witness, which has not yet been recovered.

\begin{table*}[t]
\caption{Audit status of the simplex computations.  Here $k_\Omega$ and $k_E$ are the accessible state and effect dimensions after Stage-1 reduction.  ``Dual skipped'' means that the complementary-slackness dual-recovery step was deliberately not run because the exact parameter problem was large; it is not evidence against optimality.}
\label{tab:appendix-audit}
\resizebox{\textwidth}{!}{%
\begin{tabular}{lcccccccl}
\toprule
Configuration & projectors & $k_\Omega$ & $k_E$ & facets & $r$ & exact primal & exact dual & note\\
\midrule
Firefly logic $L_{12}$ & 5 & 4 & 4 & 5 & $0$ & yes & yes & exact optimum\\
Yu--Oh 13 & 13 & 6 & 6 & 24 & $3/8$ & yes & yes & exact optimum\\
Yu--Oh 25 completion & 25 & 6 & 6 & 96 & $21/44$ & yes & no & dual skipped\\
Tkadlec Fig.~2, nonunital & 37 & 6 & 6 & 276 & $35/64$ & yes & no & dual skipped\\
Cabello 18--9 & 18 & 10 & 10 & 146 & $1/3$ & yes & no & dual skipped\\
Peres--Mermin/Peres 24 & 24 & 10 & 10 & 120 & $4/9$ & yes & no & dual skipped\\
Completed pentagon $C_5$ & 10 & 6 & 6 & 12 & $0.119140568134$ & no & no & algebraic/numerical branch\\
Specker bug/TIFS & 13 & 6 & 6 & 26 & $0.31645570$ & no & no & numerical high-precision run\\
$\Gamma_3$ & 25 & 6 & 6 & 230 & $0.50219773$ & no & no & numerical high-precision run\\
\bottomrule
\end{tabular}%
}
\end{table*}

The completed pentagon is not exact-certified by the rational backend because the coordinates in Eq.~\eqref{eq:kcbs-rays} are algebraic.  The numerical branch found 10 projectors, full accessible dimension six, 12 facets, numerical robustness $r=0.119140568134$, maximum residual approximately $1.8\times10^{-13}$, and 16 active $\sigma$ entries.  The Specker bug and $\Gamma_3$ are retained as numerical rows from the earlier high-precision run.  Exact certification of these algebraic rows would require either an algebraic-number polyhedral backend or a rationally equivalent reformulation.

\bibliography{svozil}

\end{document}

